%=====================================================================
%  CV - Jemay Salomon
%  Modele LaTeX simple, robuste et facile a maintenir (classe article)
%  Compilation : pdflatex cv_salomon.tex  (deux passes)
%=====================================================================
\documentclass[10pt,a4paper]{article}

% --- Encodage / langue ---
\usepackage[utf8]{inputenc}
\usepackage[T1]{fontenc}
\usepackage[french]{babel}
\usepackage{lmodern}

% --- Mise en page ---
\usepackage[top=1.4cm,bottom=1.4cm,left=1.7cm,right=1.7cm]{geometry}
\usepackage{xcolor}
\usepackage{titlesec}
\usepackage{enumitem}
\usepackage{tabularx}
\usepackage[expansion=false]{microtype}
\usepackage[hidelinks]{hyperref}

% --- Couleur d'accent (modifiable) ---
\definecolor{accent}{HTML}{1F4E79}   % bleu sobre

\hypersetup{
    colorlinks=true, linkcolor=accent, urlcolor=accent,
    pdftitle={CV - Jemay SALOMON}, pdfauthor={Jemay SALOMON}
}

% --- Pas d'indentation, listes compactes ---
\setlength{\parindent}{0pt}
\setlength{\emergencystretch}{2em}  % absorbe les petits debordements de justification
% Espacement par defaut des listes ; la PUCE se regle liste par liste
% via l'argument optionnel : \begin{itemize}[label=\textbullet] ... \end{itemize}
\setlist[itemize]{leftmargin=1.2em, topsep=1pt, itemsep=1pt, parsep=0pt}

% --- Style des titres de section ---
\titleformat{\section}
  {\color{accent}\large\bfseries\scshape}{}{0pt}{}[{\vspace{2pt}\color{accent}\titlerule[0.8pt]}]
\titlespacing{\section}{0pt}{10pt}{6pt}

%---------------------------------------------------------------------
% Commandes utilitaires
%---------------------------------------------------------------------
% Entree avec intitule a gauche + date a droite, puis sous-ligne (lieu/organisme)
% --- Reglages de taille (valeurs en pt : {taille}{interligne}) ---
\newcommand{\sizeTitre}{\fontsize{11}{13.2}\selectfont}  % intitule (#1)
\newcommand{\sizeDate} {\fontsize{10}{13.2}\selectfont}  % date (#4)
\newcommand{\sizeOrg}  {\fontsize{10}{12}\selectfont}    % organisme (#2)
\newcommand{\sizeLieu} {\fontsize{10}{12}\selectfont}    % lieu (#3)

% --- Commande d'entree ---
\newcommand{\cventry}[5][2pt]{%
  \par\noindent\begin{tabularx}{\textwidth}{@{}Xr@{}}
    {\sizeTitre\textbf{#2}} & {\sizeDate\textbf{#5}}\\
    {\sizeOrg #3}   & {\sizeLieu #4}\\
  \end{tabularx}\par\vspace{#1}}

% Ligne simple titre/date (sans sous-ligne)
\newcommand{\cvline}[2]{%
  \noindent\begin{tabularx}{\textwidth}{@{}Xr@{}}
    #1 & #2\\
  \end{tabularx}}

\begin{document}
\pagestyle{empty}

%=====================================================================
% EN-TETE
%=====================================================================
\begin{center}
{\Huge\bfseries\color{accent} JEMAY SALOMON}\\[2pt]
{\large Doctorant en génétique quantitative}\\[4pt]
INRAE, GQE--Le Moulon, Université Paris-Saclay, 91190 Gif-sur-Yvette, France\\[2pt]
+33 7 45 51 94 65 \quad\textbar\quad
\href{mailto:jemay.salomon@inrae.fr}{jemay.salomon@inrae.fr} \quad\textbar\quad
\href{https://jemaysalomon.github.io/}{jemaysalomon.github.io}
\end{center}

%=====================================================================
\section{Profil}
Doctorant en 3\textsuperscript{e} année en génétique quantitative. Je m'intéresse particulièrement aux approches de modélisation avancées permettant d'estimer plus précisément les effets et les (co)variances génétiques (p.\,ex. l'utilisation de la différenciation automatique couplée aux modèles linéaires mixtes en R/C++), d'évaluer la robustesse des modèles à l'aide de simulations stochastiques et, plus généralement, d'appliquer les principes de la génétique quantitative afin de mieux comprendre l'architecture génétique des caractères complexes. Je m'intéresse également à l'optimisation des schémas de sélection (évalués par simulation), ainsi qu'à l'implémentation de la sélection génomique et/ou phénomique et à l'identification du meilleur compromis entre coût de la sélection, gain génétique et maintien de la diversité génétique.

%=====================================================================
\section{Éducation}
\cventry[8pt]{Ph.D. en Génétique Quantitative}{Université Paris-Saclay, Gif-sur-Yvette, FR}{}{Oct. 2023 -- présent}
\cventry[8pt]{M.Sc. en Génétique Quantitative}{Institut Agro Montpellier, Montpellier, FR}{}{Sept. 2021 -- Sept. 2023}
\cventry{B.Sc. en Agronomie}{Université Quisqueya, Port-au-Prince, HT}{}{Sept. 2013 -- Sept. 2018}

%=====================================================================
\section{Expérience professionnelle}
\cventry[-2pt]{Assistant de recherche}{INRAE -- GQE Le Moulon, Gif-sur-Yvette, FR}{}{Oct. 2023 -- présent}
\begin{itemize}[label=\textbullet]
  \item Étudier les bases et l'architecture génétique de l'aptitude au mélange interspécifique du blé en association avec le pois.
  \item Développer une méthodologie de sélection conjointe pour les systèmes en monoculture et en association ("intercropping").
  \item Concevoir et mettre en œuvre des protocoles de recherche.
\end{itemize}
\bigskip
\cventry[-2pt]{Stagiaire (Master II)}{CHIBAS, Mirebalais, HT}{}{Mars 2023 -- Août 2023}
\begin{itemize}[label=\textbullet]
  \item Optimiser des modèles de prédiction génomique du rendement (grain) du sorgho à travers plusieurs environnements.
  \item Concevoir et mettre en œuvre des protocoles de recherche.
  \item Collecter et gérer des bases de données phénotypiques.
\end{itemize}
\bigskip
\cventry[-2pt]{Stagiaire (Master I)}{CIRAD -- CNRS, Montpellier, FR}{}{Mai 2022 -- Août 2022}
\begin{itemize}[label=\textbullet]
  \item Modéliser la relation phyllochrone--plastochrone chez le sorgho afin de prédire la date d'initiation paniculaire.
  \item Participer aux opérations d'essais au champ (collecte et analyse de données phénotypiques).
  \item Synthétiser les résultats dans un rapport scientifique (analyse des données et interprétation des modèles de croissance).
\end{itemize}

%=====================================================================
\section{Publications}
\textbf{Articles publiés}
\begin{itemize}[label=\textemdash]
  \item \textbf{Salomon, J}., Enjalbert, J. \& Flutre, T. (2026). Joint modeling of social genetic effects in mono- and pluri-specific groups: case study in intercrops. \textit{Genetics}, iyag177. \href{https://doi.org/10.1093/genetics/iyag177}{doi.org/10.1093/genetics/iyag177}
\end{itemize}
\bigskip
\textbf{Manuscrits en préparation}
\begin{itemize}[label=\textemdash, itemsep=6pt]
  \item \textbf{Salomon, J}., Levardois, K., Freitas, V., Moutier, N., Enjalbert, J. \& Flutre, T. (2026). Variability of social genetic effects in a wheat--pea intercropping trial over three environments. \textit{In prep.}
  \item Charles, J.R., \textbf{Salomon, J}. \& Pressoir, G. (2026). Genomic prediction of sorghum across environments: effects of training design and stability indices. \textit{In prep.}
\end{itemize}

\newpage
%=====================================================================
\section{Logiciels}
\begin{itemize}[label=\textemdash, itemsep=6pt]
  \item Flutre, T., Enjalbert, J., Forst, E., Remerand, M., \& \textbf{Salomon, J}. (2026). \textbf{plantmix}: Genetic Study of Plant Mixtures (v1.0.2) [Logiciel R]. \href{https://cran.r-project.org/web/packages/plantmix/index.html}{CRAN}
      \item \textbf{Salomon, J}. (2026). \textbf{TMBusers}: Useful functions to work with Template Model Builder (TMB). [Logiciel R]. \href{https://github.com/JemaySALOMON/TMBusers}{GitHub}
\end{itemize}

%=====================================================================
\section{Conférences}
\begin{itemize}[label=\textemdash, itemsep=6pt]
  \item \textbf{Salomon, J}., Levardois, K., Freitas, V., Moutier, N., Enjalbert, J. \& Flutre, T. (2026).. Dissecting the genetic architecture of direct and social genetic values: case study in wheat-pea intercrops. \textit{Présentation orale}, Mix4Agroecology. Paris, France.
  \item \textbf{Salomon, J}., Enjalbert, J., \& Flutre, T. (2025). Leveraging a broad gradient of plant--plant interactions to efficiently breed for cereal--legume mixtures. \textit{Poster}, XIX EUCARPIA Biometrics in Plant Breeding. Edinburgh,UK. \href{https://hal.inrae.fr/hal-05546246}{hal-05546246}
  \item \textbf{Salomon, J}., Enjalbert, J., \& Flutre, T. (2025). Leveraging genetic diversity to study plant--plant interactions and breed for crop mixtures. \textit{Poster}, PSL-Qlife Quantitative Biology Winter School 2025. Paris, France.
  \item \textbf{Salomon, J}., Enjalbert, J., \& Flutre, T. (2024). Leveraging a broad gradient of plant--plant interactions to efficiently breed for cereal--legume mixtures. \textit{Présentation orale}, SFE2 Conference. Lyon, France.
  \item \textbf{Salomon, J}., \& Pressoir, G. (2023). Optimization of prediction models for yield stability in sorghum: physiological determinisms, genomic prediction incorporating G$\times$E, and phenomic approaches. \textit{Poster}, ELLS Scientific Student Conference 2023. Stuttgart, Germany.
\end{itemize}

%=====================================================================
\section{Formation complémentaire}
\cventry{Atelier de statistique}{Saint-Rémy-lès-Chevreuse, FR}{}{Nov. 2025}
\noindent
\textbf{Modélisation non linéaire à effets mixtes et modélisation conjointe longitudinale--survie} : 
  \begin{itemize}[label=\textbullet, itemsep=1pt]
	 \item Modèles non linéaires à effets mixtes pour les données biologiques et longitudinales
	 \item Modélisation conjointe de données longitudinales et de survie
	 \item Sélection de variables en grande dimension à l'aide des approches LASSO et bayésiennes
	 \item Programmation scientifique reproductible et différenciation automatique
	 \item Applications de la modélisation statistique à l'amélioration des plantes avec R et Python
\end{itemize}
\bigskip
\cventry{BIOMETRICS EUCARPIA 2025 -- Satellite Workshop}{The Roslin Institute, Édimbourg, UK}{}{Sept. 2025}
\noindent
\textbf{Simulation de programmes de sélection (AlphaSimR, FieldSimR)}
  \begin{itemize}[label=\textbullet, itemsep=1pt]
  \item Simulation et optimisation de programmes de sélection à l'aide du package AlphaSimR (R)
  \item Modélisation et prise en compte de l'hétérogénéité spatiale dans les essais au champ à l'aide du package FieldSimR (R)
\end{itemize}

%=====================================================================
\section{Mentorat}
\begin{itemize}[label=\textbullet, itemsep=6pt]
  \item \textbf{Rouguiyatou Diallo} (04/2025--08/2025, co-encadrement avec T. Flutre) --- Mémoire de M1 : bases génétiques de la plasticité phénotypique du blé en culture pure versus en association avec le pois (croissance, limitation des épidémies).
  \item \textbf{Rémi Viéban} (06/2026) --- Stage 1A : évaluation de la prédiction génomique et phénomique dans le contexte des mélanges céréale--légumineuse.
\end{itemize}

%=====================================================================
\section{Compétences}
{\small
\begin{tabularx}{\textwidth}{@{}>{\raggedright\arraybackslash}X >{\raggedright\arraybackslash}X >{\raggedright\arraybackslash}X@{}}
Sélection pour les mélanges & Génétique quantitative & Amélioration des plantes\\[2pt]
Génétique statistique & Modèle linéaire mixte & Modélisation statistique\\[2pt]
Sélection génomique / phénomique & Programmation R & C++ (intermédiaire)\\
\end{tabularx}}

%=====================================================================
\section{Langues}
\begin{itemize}[label=\textbullet]
  \item Créole haïtien (bilingue ou natif)
  \item Français (bilingue ou natif)
  \item Anglais (professionnel)
\end{itemize}

%=====================================================================
\newpage
\section{Centres d'intérêt}
\begin{itemize}[label=\textbullet, itemsep=6pt]
\item Production musicale  \item DJ \& performance live.
\end{itemize}


%=====================================================================
\section{Référents}
\begin{itemize}[label=\textbullet, itemsep=6pt]
  \item \textbf{Dr Jérôme Enjalbert},  Directeur de recherche, GQE-Le Moulon (Université Paris-Saclay, INRAE, CNRS, AgroParisTech). \href{mailto:jerome.enjalbert@inrae.fr}{jerome.enjalbert@inrae.fr} : 01 69 15 68 55
  \item \textbf{Dr Timothée Flutre}, Chargé de recherche, GQE-Le Moulon (Université Paris-Saclay, INRAE, CNRS, AgroParisTech). \href{mailto:timothee.flutre@inrae.fr}{timothee.flutre@inrae.fr} : 01 69 15 61 33
  \item \textbf{Pr Dominique This}, Professeure, Institut Agro Montpellier, UMR AGAP. \href{mailto:dominique.this@supagro.fr}{dominique.this@supagro.fr} : 04 67 61 58 29
\end{itemize}

\end{document}
